% © 2026 Ing. David Jaroš — CC BY-NC-ND 4.0
%
% This work is licensed under a Creative Commons Attribution-NonCommercial-NoDerivatives
% 4.0 International License (CC BY-NC-ND 4.0).
%
% License History: Earlier drafts (up to v0.3) were released under CC BY 4.0.
% From v0.4 onward, all material is released under CC BY-NC-ND 4.0 to protect
% the integrity of the theoretical work during ongoing academic development.

\documentclass[11pt,a4paper]{article}
\usepackage{amsmath, amssymb, amsthm}
\usepackage{hyperref}
\usepackage{geometry}
\geometry{margin=2.5cm}

\newtheorem{theorem}{Theorem}
\newtheorem{proposition}{Proposition}
\newtheorem{definition}{Definition}
\newtheorem{remark}{Remark}
\newtheorem{conjecture}{Conjecture}

\title{Quantization of the \texorpdfstring{$\Theta$}{Theta}-Field in Unified Biquaternion Theory\\
\large Action, Canonical Structure, Path Integral, and Propagator}
\author{Ing.\ David Jaroš}
\date{2026}

\begin{document}
\maketitle

\begin{abstract}
We formalise the quantum field theory of the biquaternionic $\Theta$-field in
Unified Biquaternion Theory (UBT).  Starting from the canonical action
$S[\Theta]$, we identify the canonical conjugate momenta, construct the
path-integral formulation, derive the Feynman propagator, and identify the
Kaluza--Klein (KK) mode structure arising from the compact imaginary-time
direction $\psi \sim \psi + 2\pi R_\psi$.  All claims are labelled with their
proof status following the UBT notation: \textbf{[L0]} proved, \textbf{[L1]}
proved given stated assumptions, \textbf{[L2]} motivated conjecture,
\textbf{[O]} open, \textbf{[SE]} semi-empirical.
\end{abstract}

\tableofcontents

% ======================================================================
\section{Canonical \texorpdfstring{$\Theta$}{Theta}-Field Action}
\label{sec:action}
% ======================================================================

\subsection{Definition and sources}

The fundamental field of UBT is the biquaternion
$\Theta(q,\tau) \in \mathbb{H} \otimes \mathbb{C}$ defined in
\texttt{canonical/fields/theta\_field.tex} (Definition 1, status \textbf{[L0]}).
The complex time is $\tau = t + i\psi$ with $\psi \sim \psi + 2\pi R_\psi$
(Axiom A2, \texttt{THEORY/axioms/canonical\_axioms.tex}).

\subsection{Action functional}

Following \texttt{consolidation\_project/appendix\_ALPHA\_one\_loop\_biquat.tex}
(eq.\ \eqref{eq:theta-action-source}) and
\texttt{THETA\_FIELD\_DEFINITION.md} §5.1, the canonical action is:
\begin{equation}
\label{eq:action}
S[\Theta]
= \int d^4x\, dt\, d\psi\; \sqrt{|\det G|}\;
  \Bigl[
    \tfrac{1}{2} G^{MN} \langle D_M \Theta,\, D_N \Theta\rangle
    - V(\Theta)
    - \tfrac{1}{4}\langle F_{MN},\, F^{MN}\rangle
  \Bigr],
\end{equation}
where:
\begin{itemize}
  \item $G_{MN}$ is the biquaternionic metric on the $4+2$-dimensional
        manifold $(x^\mu, t, \psi)$;
  \item $D_M = \partial_M + igA_M$ is the gauge-covariant derivative;
  \item $\langle \cdot, \cdot \rangle$ is the Hermitian inner product on
        $\mathbb{H}\otimes\mathbb{C}$;
  \item $V(\Theta)$ is the Higgs-type potential whose form is given in
        \texttt{canonical/interactions/higgs.tex};
  \item $F_{MN} = \partial_M A_N - \partial_N A_M - ig[A_M, A_N]$ is the
        gauge field strength.
\end{itemize}

\begin{remark}[Status]
  The action \eqref{eq:action} is the canonical UBT action derived from the
  Hermitian structure of $\mathbb{H}\otimes\mathbb{C}$.  Its full derivation,
  including boundary terms, is in
  \texttt{consolidation\_project/appendix\_A\_theta\_action.tex}.
  Status: \textbf{[L1]} — proved given the axioms.
\end{remark}

% ======================================================================
\section{Canonical Conjugate Variables}
\label{sec:canonical}
% ======================================================================

\subsection{Separation into physical and imaginary sectors}

In the isotropic limit $G_{MN} = \mathrm{diag}(g_{\mu\nu},-1,-1)$, the
action splits as
\begin{equation}
  S = S_{\mathrm{4D}} + S_\psi,
\end{equation}
where $S_{\mathrm{4D}}$ is the four-dimensional part and $S_\psi$ is the
kinetic term in the $\psi$-direction.

\subsection{Canonical momenta}

For a generic component $\Theta_a$ ($a = 0,1,2,3$) in the quaternion basis
$\{1, i, j, k\}$, the Euler--Lagrange formalism yields the canonical momenta:
\begin{equation}
\label{eq:pi}
  \Pi^a(x, t, \psi) \;=\; \frac{\partial \mathcal{L}}{\partial(\partial_t \Theta_a)}
  \;=\; \dot{\Theta}_a(x, t, \psi),
\end{equation}
where a dot denotes $\partial_t$.

\begin{definition}[Equal-time commutation relations]\label{def:etcr}
After canonical quantisation:
\begin{equation}
\label{eq:etcr}
  \bigl[\Theta_a(x, t, \psi),\, \Pi^b(x', t, \psi')\bigr]
  = i\hbar\, \delta^{ab}\, \delta^{(3)}(x - x')\, \delta(\psi - \psi').
\end{equation}
\end{definition}

\begin{remark}[Status]
  The commutation relation \eqref{eq:etcr} is a direct application of
  standard canonical quantisation to the biquaternionic field.
  Status: \textbf{[L1]} — proved given action \eqref{eq:action}.
  The non-trivial content lies in the $\delta(\psi-\psi')$ factor which
  implies that modes at different $\psi$-slices are independent quantum
  degrees of freedom.
\end{remark}

% ======================================================================
\section{Path-Integral Formulation}
\label{sec:path_integral}
% ======================================================================

\subsection{Partition function}

The generating functional is:
\begin{equation}
\label{eq:Z}
  Z[J]
  = \int \mathcal{D}\Theta\; \exp\!\Bigl(
      \frac{i}{\hbar} S[\Theta]
      + \frac{i}{\hbar}\int J_a \Theta^a
    \Bigr),
\end{equation}
where $J_a(x,\tau)$ are biquaternionic source fields.

\subsection{Gaussian sector and free propagator}

In the free (non-interacting) limit, $V(\Theta) = \tfrac{1}{2}m^2\langle\Theta,\Theta\rangle$
and the gauge fields are set to zero.  The action becomes:
\begin{equation}
\label{eq:S_free}
  S_{\mathrm{free}}[\Theta]
  = \int d^4x\, d\tau\; \Bigl[
    \tfrac{1}{2} G^{MN} \partial_M \Theta^\dagger \partial_N \Theta
    - \tfrac{1}{2} m^2 |\Theta|^2
  \Bigr].
\end{equation}

After Fourier transform in all directions, including the KK expansion on the
compact $\psi$-circle (see Section~\ref{sec:kk}), one obtains the free
inverse propagator for KK mode $n$:
\begin{equation}
\label{eq:inv_prop}
  \tilde{G}^{-1}(k, n) \;=\; k^\mu k_\mu + \frac{n^2}{R_\psi^2} + m^2,
\end{equation}
where $k^\mu$ is the 4-momentum and $n \in \mathbb{Z}$ is the KK mode number.

\begin{definition}[Feynman Propagator]\label{def:propagator}
The Feynman propagator for KK mode $n$ in momentum space is:
\begin{equation}
\label{eq:propagator}
  \boxed{
  \Delta_n(k) \;=\; \frac{i}{k^\mu k_\mu - m_n^2 + i\varepsilon},
  \qquad
  m_n^2 \;=\; m^2 + \frac{n^2}{R_\psi^2}.
  }
\end{equation}
\end{definition}

\begin{remark}[Biquaternionic tensor structure]
  The full propagator carries tensor indices in the biquaternionic algebra
  $\mathbb{H}\otimes\mathbb{C}$.  In the component basis it becomes a
  $4\times 4$ complex matrix in quaternionic index space.  For a single
  component $\Theta_a$, Definition~\ref{def:propagator} applies component-wise.
  Status: \textbf{[L1]}.
\end{remark}

% ======================================================================
\section{Kaluza--Klein Mode Expansion}
\label{sec:kk}
% ======================================================================

\subsection{Fourier expansion on the \texorpdfstring{$\psi$}{psi}-circle}

Following \texttt{report/theta\_spectrum\_analysis.md} §2.2 and
\texttt{canonical/geometry/phase\_projection.tex} §KK, the compactification
$\psi \sim \psi + 2\pi R_\psi$ yields the mode expansion:
\begin{equation}
\label{eq:kk_expansion}
  \Theta(x, t, \psi)
  = \sum_{n \in \mathbb{Z}} \hat\Theta_n(x, t)\;
    e^{i n \psi / R_\psi}.
\end{equation}

\begin{proposition}[KK mass spectrum]\label{prop:kk_masses}
  The 4D effective mass of mode $n$ is:
  \begin{equation}
  \label{eq:kk_masses}
    m_n^2 = m^2 + \frac{n^2}{R_\psi^2}.
  \end{equation}
  The $n = 0$ mode is exactly massless (in the $m = 0$ limit) by translation
  invariance in $\psi$.
  \textbf{Status: [L0]} — algebraic identity.
\end{proposition}

\subsection{Mode commutation relations}

Inserting \eqref{eq:kk_expansion} into \eqref{eq:etcr} and using
orthonormality $\int_0^{2\pi R_\psi} d\psi\; e^{i(n-m)\psi/R_\psi}
= 2\pi R_\psi\, \delta_{nm}$:
\begin{equation}
\label{eq:mode_etcr}
  \bigl[\hat\Theta_{n,a}(x),\; \hat\Pi^{m,b}(x')\bigr]
  = \frac{i\hbar}{2\pi R_\psi}\, \delta_{nm}\, \delta^{ab}\,
    \delta^{(3)}(x - x').
\end{equation}

\begin{remark}[Status]
  The mode algebra \eqref{eq:mode_etcr} is a straightforward Fourier decomposition
  of the canonical commutators.  Status: \textbf{[L1]}.
\end{remark}

\subsection{Fermionic representations in the mode algebra}

The question of whether any modes form fermionic representations is
addressed in \texttt{research/theta\_fermion\_emergence.tex}.  Here we note
only the structural observation:

\begin{remark}[Winding parity]
  The $\psi$-winding number $n$ provides a $\mathbb{Z}_2$ parity by
  $(-1)^n$.  Modes with odd $n$ are $\psi$-odd.  This grading is
  the algebraic seed of the fermionic sector; see
  \texttt{research/theta\_fermion\_emergence.tex} for the detailed analysis.
\end{remark}

% ======================================================================
\section{Interaction Vertices and Feynman Rules}
\label{sec:vertices}
% ======================================================================

\subsection{Cubic and quartic vertices from the potential}

Expanding the potential $V(\Theta)$ around the vacuum expectation value
$\langle\Theta\rangle = v$ (Higgs mechanism, canonical source:
\texttt{canonical/interactions/higgs.tex}):
\begin{equation}
  V(\Theta) = \tfrac{1}{4}\lambda(|\Theta|^2 - v^2)^2
  = \lambda v^2 |\delta\Theta|^2 + \lambda v |\delta\Theta|^3
    + \tfrac{1}{4}\lambda |\delta\Theta|^4 + \text{const.},
\end{equation}
where $\delta\Theta = \Theta - v$.  This produces the standard cubic (Yukawa-type)
and quartic self-interaction vertices.

\subsection{KK mode vertices}

KK number is conserved at each vertex due to the periodicity of $e^{in\psi/R_\psi}$:
\begin{equation}
\label{eq:kk_conservation}
  n_1 + n_2 + \cdots + n_k = 0 \pmod{1}
  \quad \text{(at each $k$-point vertex).}
\end{equation}

This is the KK mode-number conservation law (\textbf{[L0]}).

% ======================================================================
\section{One-Loop Effective Action}
\label{sec:one_loop}
% ======================================================================

\subsection{Vacuum polarization from KK modes}

Each KK mode $n \neq 0$ contributes to the vacuum polarization through the
standard one-loop diagram.  The leading contribution from mode $n$ is
\begin{equation}
  \Pi^{(n)}(k^2)
  = \frac{g^2}{16\pi^2} \int_0^1 dx\; x(1-x)\,
    \ln\!\Bigl(\frac{m_n^2 - x(1-x)k^2}{\mu^2}\Bigr),
\end{equation}
following the detailed calculation in
\texttt{consolidation\_project/appendix\_ALPHA\_one\_loop\_biquat.tex} §B.3.

\begin{remark}[Connection to B\_base]
  The KK sum $\sum_{n\in\mathbb{Z}} \Pi^{(n)}$ is directly related to
  the $B_{\mathrm{base}}$ coefficient appearing in the fine structure constant
  formula.  The current status of this connection is documented in
  \texttt{research/B\_base\_spectral\_determinant.tex} and
  \texttt{canonical/bridges/QED\_limit\_bridge.tex}.
\end{remark}

% ======================================================================
\section{Open Problems and Research Tracks}
\label{sec:open}
% ======================================================================

\begin{enumerate}
  \item \textbf{[O] Off-shell closure}: The Θ-field equation
    $\nabla^\dagger\nabla\Theta = \kappa\mathcal{T}$ is verified on-shell
    (Steps 1--5 of \texttt{canonical/bridges/GR\_chain\_bridge.tex}).
    Off-shell closure is an identified open gap.

  \item \textbf{[O] R\_psi from first principles}: The compactification
    radius $R_\psi$ is a free parameter.  Fixing it dynamically (e.g.,
    from a moduli stabilisation mechanism) is open; see
    \texttt{AUDITS/complex\_time\_modular\_audit.md}.

  \item \textbf{[O] Fermionic statistics from KK algebra}:
    Whether odd-winding modes obey fermionic commutation relations
    (anticommutators) requires showing that the biquaternionic algebra
    supports a consistent $\mathbb{Z}_2$ grading at the quantum level.
    See \texttt{research/theta\_fermion\_emergence.tex}.

  \item \textbf{[O] Unitarity and renormalisability}:
    The infinite KK tower raises questions of unitarity and
    renormalisability.  Standard decoupling arguments suggest the
    tower is safe for $R_\psi^{-1} \gg \mu$, but a rigorous proof
    in the biquaternionic context is open.
\end{enumerate}

% ======================================================================
\section{Cross-References}
\label{sec:xref}
% ======================================================================

\begin{tabular}{ll}
\hline
Topic & File \\
\hline
Canonical Θ-field definition & \texttt{canonical/fields/theta\_field.tex} \\
Θ-field action (full derivation) & \texttt{consolidation\_project/appendix\_A\_theta\_action.tex} \\
KK spectrum analysis & \texttt{report/theta\_spectrum\_analysis.md} \\
KK phase projection & \texttt{canonical/geometry/phase\_projection.tex} \\
Fermion emergence & \texttt{research/theta\_fermion\_emergence.tex} \\
Bridge: quantum structure & \texttt{canonical/bridges/theta\_quantum\_structure.tex} \\
One-loop $\alpha$ calculation & \texttt{consolidation\_project/appendix\_ALPHA\_one\_loop\_biquat.tex} \\
QED limit chain & \texttt{canonical/bridges/QED\_complete\_chain.tex} \\
\hline
\end{tabular}

\end{document}
